\subsection{Политическая культура в контексте концепций политического порядка и модернизации С. Хантингтона}

Важнейшим дополнением к классическим теориям политической культуры стал институциональный подход, развитый американским политологом С. Хантингтоном в его фундаментальном труде «Политический порядок в меняющихся обществах». В отличие от Г. Алмонда и С. Вербы, рассматривавших политическую культуру преимущественно сквозь призму психологической стабильности западных демократий, С. Хантингтон исследовал данный феномен в динамическом контексте глобальных процессов модернизации и их влияния на устойчивость политических систем различного типа\footnote{Хантингтон С. Политический порядок в меняющихся обществах. --- М.: Прогресс-Традиция, 2004. --- С. 19.}.

В рамках концепции С. Хантингтона ключевой детерминантой стабильности общества выступает соотношение между уровнем политического участия масс и степенью институционализации самой политической системы. Модернизация неизбежно влечет за собой трансформацию традиционной политической культуры: происходят глубинные изменения в сознании людей, урбанизация, быстрый рост грамотности и средств массовой коммуникации. Всё это приводит к резкому повышению политической активности населения и расширению его требований к институтам власти. Однако если каналы интеграции этой активности (политические партии, избирательные и правовые механизмы) развиты слабо, возникает ценностно-нормативный дисбаланс, ведущий к политическому упадку и дестабилизации порядка.

В зависимости от характера нормативно-поведенческих паттернов взаимодействия общества и государства С. Хантингтон выделяет два ключевых типа систем, основанных на специфических типах политической культуры:

1. Гражданские системы, характеризующиеся высокой степенью институционализации. В таких обществах доминирует развитая правовая политическая культура, в рамках которой различные социальные группы соглашаются использовать исключительно легитимные, институционализированные каналы для выражения своих интересов. Политическая элита обладает высокой автономией, а граждане разделяют ценности процедурного консенсуса, что позволяет мирно разрешать возникающие социальные конфликты.

2. Преторианские системы, отличающиеся низким уровнем институционализации при высоком уровне политической активности масс. В преторианской политической культуре отсутствует общее согласие относительно легитимных методов разрешения споров\footnote{Хантингтон С. Политический порядок в меняющихся обществах. --- М.: Прогресс-Традиция, 2004. --- С. 106.}. Различные социальные группы (армия, студенчество, духовенство, профсоюзы) действуют в политическом поле напрямую, используя собственные специфические средства: военные перевороты, забастовки, погромы и прямое насилие. Политическая культура здесь фрагментирована, а институты власти воспринимаются не как суверенные арбитры, а как инструменты борьбы отдельных группировок.

Следовательно, с точки зрения С. Хантингтона, успешное становление устойчивой политической культуры требует не просто декларации абстрактных свобод, а долгосрочного формирования прочных политических организаций, способных абсорбировать и упорядочивать возрастающую активность масс\footnote{Хантингтон С. Политический порядок в меняющихся обществах. --- М.: Прогресс-Традиция, 2004. --- С. 54.}. Политическая модернизация признается успешной лишь тогда, когда темпы развития институциональной культуры опережают темпы мобилизации населения, создавая тем самым прочный политический порядок.